% ============================================================
% pretty-charts TikZ/pgfplots 导言区模板 —— academic 档（T1 出版级）
% 编译方式：XeLaTeX。business/showcase 档改字号与 cycle list（见文件末尾注释）
% ============================================================
\usepackage{pgfplots}
\pgfplotsset{compat=1.18}
\usetikzlibrary{arrows.meta, positioning, calc}

% ---- 图表字体：中文黑体 + 西文 Arial（XeLaTeX；正文仍可用全局宋体+TNR，互不影响）----
% 注意：Windows 的 Noto Sans SC 多为可变字体(VF)，xdvipdfmx 无法嵌入（fatal: Invalid font），
% 故 TikZ 导出 PDF 用 Microsoft YaHei；装有静态版思源黑体(Source Han Sans SC)的机器可换回。
\usepackage{fontspec}
\usepackage{xeCJK}
\newCJKfontfamily\chartcjk{Microsoft YaHei}
\newcommand{\chartfont}[1]{{\chartcjk\fontspec{Arial}#1}}

% ---- 色板：academic（Okabe-Ito 重排，与 references/style/palettes/academic.json 一致）----
\definecolor{pcBlue}{HTML}{0072B2}
\definecolor{pcOrange}{HTML}{D55E00}
\definecolor{pcGreen}{HTML}{009E73}
\definecolor{pcMagenta}{HTML}{CC79A7}
\definecolor{pcYellow}{HTML}{E69F00}
\definecolor{pcCyan}{HTML}{56B4E9}
\definecolor{pcGray}{HTML}{999999}
\definecolor{pcGrid}{HTML}{CCCCCC}
\definecolor{pcAxis}{HTML}{333333}
\definecolor{pcTick}{HTML}{4D4D4D}
\definecolor{pcTitleBlack}{HTML}{1A1A1A} % = palettes/academic.json text.title
\definecolor{pcNeutral}{HTML}{B3B3B3}    % = palettes/academic.json neutral（中性灰，style-guide §3.2.4）

% ---- 全局轴风格：去顶右刺、浅色 y 网格、字号层级 ----
\pgfplotsset{
  every axis/.append style={
    axis line style={pcAxis, line width=0.4pt},
    label style={font=\small, text=pcAxis},
    tick label style={font=\footnotesize, text=pcTick},
    title style={font=\small\bfseries, text=pcTitleBlack},
    legend style={font=\footnotesize, draw=none, fill=none, nodes={scale=0.9, transform shape}},
    grid=major,
    major grid style={pcGrid, line width=0.4pt},
    axis x line*=bottom,
    axis y line*=left,
    tick style={pcTick, line width=0.4pt},
  },
  every axis plot/.append style={line width=0.9pt},
  /pgfplots/error bars/error bar style={line width=0.5pt},
}

% ---- 色板循环列表 ----
\pgfplotscreateplotcyclelist{pc academic}{
  pcBlue\\ pcOrange\\ pcGreen\\ pcMagenta\\ pcYellow\\ pcCyan\\ pcGray\\
}

% ---- 用法示例 ----
% \begin{tikzpicture}
%   \begin{axis}[cycle list name=pc academic, xlabel={年份}, ylabel={销量}, title={标题}]
%     \addplot coordinates {(1,2) (2,4) (3,3)};
%   \end{axis}
% \end{tikzpicture}

% ============================================================
% business 档（T2）差异：字号放大一档（label \normalsize / tick \small），
% cycle list 换 4E79A7/F28E2B/E15759/76B7B2/59A14F/EDC948/B07AA1/FF9DA7/9C755F/BAB0AC
% showcase 档（T3）：字号再放大（label \large / tick \normalsize），line width=1.4pt，
% cycle list 换 0077BB/EE7733/009988/EE3377/33BBEE/CC3311/BBBBBB
% ============================================================
